% !TEX root = ../main.tex

\section{METHODOLOGY}
\label{sec:method}

% =====================================================================
% Chapter 3 rewritten 18 Nov 2026 around the risk-calibrated payload
% router; the per-utterance stopping-time scheduler of the concept
% defense is now preliminary work (Appendix D).
% 28 Nov 2026: written as a plan; every candidate path is calibrated and
% evaluated separately and the paths are compared in the evaluation (the best
% path is a finding, not part of the architecture); Figure 3.2, Remark 3.1 and
% computed values removed; Section 3.1 shortened.
% 27 Nov 2026: candidate teacher paths (k = 2, 4, 6 steps on the 8-interval
% grid; 10, 13, 16 passes). Learn-Then-Test selects the least costly path and
% its threshold (fixed-sequence test per path at delta/3, union bound), in line
% with the 27 Nov Chapter 1. The 2- and 6-step paths are no longer ablations.
% Rules kept in this draft: future tense for what the study will do,
% present tense for facts about released code and for mathematics; no
% result values (outputs are symbols or "will be reported"); every
% display equation has a "where:" list and a List of Equations entry;
% objectives appear only in headings and in the workflow (Table 3.2).
% =====================================================================

% =====================================================================
\subsection{Research Design}
\label{sec:design}

This study will use a quantitative experimental design combined with algorithm development and statistical calibration. It will study DMOSpeech~2 \parencite{li2025dmospeech2}, a zero-shot text-to-speech system in which a large teacher, F5-TTS \parencite{chen2024f5}, carries out the first generation steps and a smaller, distilled student finishes the speech. Both networks work on a mel-spectrogram, a picture of the sound over time and frequency, which a vocoder turns into audio. The proposed method, a payload router, will let the teacher take fewer, larger steps and will decide for each sentence whether the teacher hands the student its own estimate of the final mel-spectrogram or a corrected version of it. Three candidate teacher paths, which differ in the number of teacher steps, will be studied. The method therefore adapts what is handed over for each sentence, not how much computation a sentence receives. Computing cost will be counted in \emph{forward passes}, the number of times a network is run.

The study will follow the three specific objectives (Section~\ref{sec:workflow}). It will first reproduce the released system and describe the two payloads on development speakers. It will then formulate the router and, with Learn-Then-Test, calibrate its threshold for each candidate path on separate calibration speakers. Finally, it will evaluate every calibrated path against the established systems on unseen test speakers and a second dataset; comparing the paths will show which one gives the lowest cost while keeping excess failures within the stated limit.

All comparisons will be paired: every system will synthesize the same inputs from the same random starting values, and no network will be retrained. The prompt condition and the number of calibration speakers will be varied; the networks, the candidate paths, the seeds and the hardware will be held fixed. The unit of analysis will be the speaker, with one sentence per speaker in the calibration and test groups. Results on the development speakers, including those of Appendix~\ref{app:snr-r}, will be treated as exploratory. The study will measure speech quality, prosodic variety, computing cost and risk, not clinical or user outcomes.

% =====================================================================
\subsection{System Architecture}
\label{sec:architecture}

The proposed method adds a decision step, the \emph{router}, to DMOSpeech~2 without changing any network. Figure~\ref{fig:method} shows the architecture, and Table~\ref{tab:components} lists its components, which fall into three groups: the frozen backbone, the released networks used without retraining; the router, which decides during generation what the teacher hands to the student; and the offline calibrator, which sets the router's threshold for each candidate teacher path once, before the system is used.

The frozen backbone consists of the duration predictor, the teacher, the distilled student and the vocoder. The duration predictor fixes the length of the output. The teacher starts from random noise and turns it step by step into a mel-spectrogram. It will take $k$ large steps, where $k$ is one of the candidate paths (two, four or six steps), then stop at the point $t_k$ on its time scale and give its estimate of the final mel-spectrogram (Section~\ref{sec:baseline}). The student runs once from pure noise to produce its draft, and $4-k/2$ more times from the selected payload to finish generation. The vocoder turns the final mel-spectrogram into a waveform.

The router sits between the teacher and the student. It reads three things: the prompt audio, from which it estimates the signal-to-noise ratio $\hat S$, a measure of how loud the speech is compared with the background noise (Definition~\ref{def:snr}); the teacher's estimate $\hat x_1$; and the student draft $s_0$. From the last two it forms the corrected payload, called the quantile projection $Q(\hat x_1,s_0)$ (Definition~\ref{def:quantile}). It also computes a score called the invariant residual $r$ (Definition~\ref{def:residual}), which measures how far the teacher's loudness values have moved from the draft's, for example because the teacher copied the prompt's background noise. The routing rule (Definition~\ref{def:policy}) passes the corrected payload when the prompt seems noisy or when the residual is above the threshold, and the teacher's own estimate otherwise. The router runs no network of its own: the draft is the only added forward pass, and it is counted in the cost (Proposition~\ref{prop:cost}).

The calibrator runs once, before deployment. For each candidate teacher path, it replays the routing rule on recorded outputs of a separate group of calibration speakers, trying a list of candidate thresholds. It keeps the largest threshold that passes a statistical test on the number of failures, which is the one that lets the most teacher estimates through (Algorithm~\ref{alg:ltt}). It changes no network and adds nothing to generation; its only output is one threshold $\hat\lambda_k$ for each path. A path for which no threshold passes falls back to the released DMOSpeech~2 system (Proposition~\ref{prop:fallback}).

% Figure (Chapter 3 v2): architecture of the risk-calibrated payload router.
% Grey boxes: released DMOSpeech 2 / F5-TTS networks (frozen). Dashed blue
% boxes: components added by the method. Every inference-time box shows its
% forward-pass cost; the total is 13 passes for every utterance.
\begin{figure}[!htbp]
\centering
\begin{adjustbox}{max width=\linewidth}
\begin{tikzpicture}[
  font=\figfont,
  base/.style={draw, rounded corners=2pt, fill=gray!8, text width=2.35cm, align=center, minimum height=1.3cm, inner sep=3pt},
  prop/.style={draw=blue!60!black, thick, dashed, rounded corners=2pt, fill=blue!5, text width=2.35cm, align=center, minimum height=1.3cm, inner sep=3pt},
  dec/.style={diamond, draw=blue!60!black, thick, dashed, fill=blue!5, aspect=1.6, inner sep=0pt, align=center, text width=1.55cm},
  arr/.style={-{Stealth[length=2mm]}, thick}
]
% inputs and frozen front end
\node[base] (in)    at (0,0)       {Prompt audio,\\prompt text,\\target text};
\node[base] (dur)   at (0,-2.6)    {Duration\\predictor\\(as released)};
\node[base] (teach) at (3.6,0)     {Teacher: 4 guided\\Euler steps to\\$t_4=0.2929$,\\endpoint $\hat x_1$\\(10 passes)};
\node[base] (draft) at (3.6,-2.6)  {Student draft\\$s_0$ from noise\\(1 pass)};
% added inference-time components
\node[prop] (snr)   at (3.6,2.5)   {Prompt SNR\\estimate $\hat S$\\(no pass)};
\node[prop, text width=2.9cm] (qr) at (7.7,-1.7) {Quantile projection\\$Q(\hat x_1,s_0)$ and\\invariant residual $r$\\(no pass)};
\node[dec]  (route) at (11.0,0)    {$\hat S<\theta$\\or\\$r>\hat\lambda_r$?};
% frozen back end
\node[base] (stud)  at (15.4,0)    {Student: 2 steps\\from the payload,\\$j=2$ (2 passes)};
\node[base] (voc)   at (15.4,-2.6) {Vocos vocoder\\$\rightarrow$ waveform};
% offline calibrator
\node[prop] (cal)   at (15.4,2.5)  {Offline calibrator\\(Learn-Then-Test)\\$\rightarrow\hat\lambda_r$};

\draw[arr] (in) -- (teach);
\draw[arr] (in) -- (dur);
\draw[arr] (dur) -- (draft);
\draw[arr] (in.north) |- (snr.west);
\draw[arr] (teach.east) -- node[above, font=\figsmall, pos=0.55] {$\hat x_1$} (route.west);
\draw[arr] ($(teach.east)+(0.4,0)$) |- ($(qr.west)+(0,0.3)$);
\draw[arr] (draft.east) -- ++(0.75,0) |- ($(qr.west)+(0,-0.3)$);
\draw[arr] (qr.east) -| node[pos=0.3, below, font=\figsmall] {$Q$, $r$} (route.south);
\draw[arr] (snr.east) -| node[pos=0.3, above, font=\figsmall] {$\hat S$} ($(route.north)+(-0.3,0)$);
\draw[arr] (cal.west) -| node[pos=0.3, above, font=\figsmall] {$\hat\lambda_r$} ($(route.north)+(0.3,0)$);
\draw[arr] (route.east) -- (stud.west);
\node[font=\figsmall, align=center] at ($(route.east)!0.5!(stud.west)+(0,0.42)$) {yes: $Q$};
\node[font=\figsmall, align=center] at ($(route.east)!0.5!(stud.west)+(0,-0.42)$) {no: $\hat x_1$};
\draw[arr] (stud) -- (voc);

% legend and cost line
\node[base, minimum height=0.35cm, text width=0.45cm] (l1) at (-0.75,-4.3) {};
\node[right, font=\figsmall] at (l1.east) {released network, weights frozen};
\node[prop, minimum height=0.35cm, text width=0.45cm] (l2) at (5.0,-4.3) {};
\node[right, font=\figsmall] at (l2.east) {added by the proposed method};
\node[right, font=\figsmall] at (10.6,-4.3) {Cost: $10+1+2=13$ passes for every input};
\end{tikzpicture}
\end{adjustbox}
\caption{Architecture of the proposed risk-calibrated payload router. The teacher always stops at the same coarse commit point; the router decides only \emph{what} the teacher hands to the student: its own endpoint estimate $\hat x_1$, or the quantile projection $Q(\hat x_1,s_0)$, which keeps the teacher's timing and uses the student draft's spectral values. The single threshold $\hat\lambda_r$ is set offline by Learn-Then-Test.}
\label{fig:method}
\end{figure}


\begin{table}[!htbp]
\centering\small
\caption{Components of the proposed architecture}
\label{tab:components}
\begin{tabularx}{\textwidth}{@{}>{\raggedright\arraybackslash}p{3.1cm} >{\raggedright\arraybackslash}X >{\raggedright\arraybackslash}p{2.3cm} >{\raggedright\arraybackslash}p{2.5cm}@{}}
\toprule
\textbf{Component} & \textbf{Role} & \textbf{Status} & \textbf{Cost per utterance} \\
\midrule
Duration predictor & Sets the number of frames to generate & Frozen & As released \\
Teacher (F5-TTS) & $k\in\{2,4,6\}$ guided Euler steps on the 8-interval grid and the endpoint estimate at $t_k$ & Frozen & $2k+2$ passes \\
Student (DMOSpeech~2), draft & One step from pure noise, giving $s_0$ & Frozen & 1 pass \\
Student (DMOSpeech~2), finish & Re-noising steps from the payload, from $j=k/2$ & Frozen & $4-k/2$ passes \\
Vocoder (Vocos) & Mel-spectrogram to waveform & Frozen & As released \\
Prompt SNR estimator & Frame-energy estimate $\hat S$ & Added, inference & No pass \\
Quantile projection & Payload $Q(\hat x_1,s_0)$ & Added, inference & No pass \\
Residual scorer & Invariant residual $r(\hat x_1,s_0)$ & Added, inference & No pass \\
Payload router $\pi_{k,\hat\lambda_k}$ & Chooses $\hat x_1$ or $Q(\hat x_1,s_0)$ & Added, inference & No pass \\
Calibrator & Sets $\hat\lambda_k$ for each path with a guarantee & Added, offline & Run once \\
\bottomrule
\end{tabularx}
\end{table}

Several choices in this architecture rest on the team's earlier experiments. Appendix~\ref{app:prelim} reports the stopping-time scheduler of the earlier proposal and the reasons it was replaced: its signal for deciding when to stop could not be shown to save computation safely, and its noise check missed some kinds of damaged prompts. Appendix~\ref{app:snr-r} reports the development runs of the router itself. These runs used few speakers and a threshold chosen without a statistical test, so they inform the design but are not evidence for the method.

% =====================================================================
\subsection{Research Workflow}
\label{sec:workflow}

Table~\ref{tab:steps} lists the fifteen steps of the methodology, grouped into the common setup and one stage per specific objective, with the section where each step is explained.

{\small
\setlength{\LTpre}{6pt}\setlength{\LTpost}{6pt}
\begin{longtable}{@{}>{\centering\arraybackslash}p{0.8cm} >{\raggedright\arraybackslash}p{8.0cm} >{\raggedright\arraybackslash}p{2.4cm} >{\raggedright\arraybackslash}p{1.5cm}@{}}
\caption{Steps of the methodology}
\label{tab:steps} \\
\toprule
\textbf{Step} & \textbf{What will be done} & \textbf{Data} & \textbf{Section} \\
\midrule
\endfirsthead
\multicolumn{4}{@{}l}{\textit{Table~\ref{tab:steps} (continued)}} \\
\toprule
\textbf{Step} & \textbf{What will be done} & \textbf{Data} & \textbf{Section} \\
\midrule
\endhead
\midrule
\multicolumn{4}{r@{}}{\textit{Continued on next page}} \\
\endfoot
\bottomrule
\endlastfoot
\multicolumn{4}{@{}l}{\textit{Common setup}} \\
1 & Divide the Seed-TTS speakers into separate development, calibration and test groups, with no speaker in more than one group and one sentence per speaker; select the LibriSpeech-PC test sentences. & All splits & \ref{sec:data} \\
2 & Create the clean prompt and seven degraded versions of every prompt; assign one condition per calibration and test speaker for the calibration and the bound check. & All splits & \ref{sec:data} \\
3 & Check that the same input always gives the same output, that the payload computations behave as defined, that the counted forward passes match Proposition~\ref{prop:cost}, and that the fallback reproduces the released system. & --- & \ref{sec:tools} \\
\midrule
\multicolumn{4}{@{}l}{\textit{Objective 1: characterize the baseline and the handover payloads}} \\
4 & Run the established systems (Table~\ref{tab:systems}) and both payloads of every candidate teacher path under every prompt condition; record WER, SIM, UTMOS, forward passes and runtime. & Development & \ref{sec:obj1} \\
5 & Run the released DMOSpeech~2 system on the full Seed-TTS test-en set with clean prompts and compare its scores with the published ones. & Seed-TTS test-en & \ref{sec:obj1} \\
6 & Fix the minimum SIM and UTMOS levels of the failure rule from the reproduced baseline; record the values of $\hat S$ and $r$. & Development & \ref{sec:obj1} \\
\midrule
\multicolumn{4}{@{}l}{\textit{Objective 2: develop the method (formulation and analysis)}} \\
7 & State the problem, define the quantile projection, the invariant residual and the routing rule, and prove their properties. & --- & \ref{sec:obj2} \\
\midrule
\multicolumn{4}{@{}l}{\textit{Objective 2: develop the method (calibration with a reliability guarantee)}} \\
8 & Fix $\alpha$, $\delta$, the candidate teacher paths, the ordered grid of $\lambda$ and the failure rule before calibration. & --- & \ref{sec:calibration} \\
9 & Record the outputs of both payloads of every path and of the released system for the calibration speakers; for each path, test the candidate thresholds in a fixed order with a statistical test on their numbers of failures (Algorithm~\ref{alg:ltt}). & Calibration & \ref{sec:calibration} \\
10 & Lock the accepted threshold $\hat\lambda_k$ of every path; a path with no accepted threshold falls back to the released system. & Calibration & \ref{sec:calibration} \\
\midrule
\multicolumn{4}{@{}l}{\textit{Objective 3: evaluate the method and test its robustness}} \\
11 & Measure on the test speakers how often the method fails where the released system does not, with a confidence interval; repeat the calibration on random re-divisions of the speakers. & Calibration, test & \ref{sec:verify} \\
12 & Run all compared systems (Table~\ref{tab:systems}), including the router on every candidate path, on the test split and on LibriSpeech-PC. & Test, LibriSpeech-PC & \ref{sec:compared} \\
13 & Compare forward passes, running time on the Tesla T4, WER, SIM, UTMOS, pitch variety and failure rates with paired statistical tests and tests of equivalence, and compare the candidate paths with one another. & Test, LibriSpeech-PC & \ref{sec:stats} \\
14 & Report results per prompt condition; calibrate on clean prompts and test on damaged ones; repeat calibration with fewer speakers and with a larger tolerance. & Calibration, test & \ref{sec:robust} \\
\midrule
15 & Report all results, including negative ones. & --- & \ref{sec:workflow} \\
\end{longtable}
}


% =====================================================================
\subsection{Common Experimental Setup}
\label{sec:setup}

\subsubsection{Baseline Sampler and Cost}
\label{sec:baseline}

The study will use the trained DMOSpeech~2 networks and generation code released with \textcite{li2025dmospeech2}. For a prompt and a target text, the duration predictor fixes the total length. The mel-spectrogram $x\in\R^{F\times 100}$ holds $F$ time frames of 100 frequency bands each: the frames of the prompt followed by the frames to be generated, which a mask $M$ marks. This section describes the released generation procedure, which is the baseline, and then the shorter teacher paths of the router.

\phead{Released sampler} The backbone is used exactly as released, so it is described here only as far as the method depends on it; Algorithm~\ref{alg:baseline} gives the complete procedure. The teacher is a flow-matching model: it starts from random noise at time $t=0$ and moves it in small steps towards speech at time $t=1$. The steps follow the sway schedule of F5-TTS \parencite{chen2024f5}, which places more steps near the start, where the coarse structure of the speech is decided:
\begin{equation}
t_i = 1-\cos\!\Big(\frac{\pi}{2}\cdot\frac{i}{N}\Big), \qquad i = 0,1,\dots,N.
\label{eq:grid}
\end{equation}
\eqentry{Teacher time grid}
\begin{whereterms}[2.6cm]
$t_i$ & time of teacher grid point $i$ on the flow axis, from $t_0=0$ (noise) to $t_N=1$ (data) \\
$N$ & number of grid intervals: $N=16$ in the released system, $N=8$ for the router
\end{whereterms}
At each grid point the teacher computes a direction of change, the guided velocity $v$, and moves its current state $x_{t_i}$ a short distance in that direction to the next grid point (an Euler step). Computing $v$ takes two forward passes, one with the text and prompt and one without, because the difference between the two is used to make the output follow the text and the voice more closely (classifier-free guidance, strength $\gamma=2$). The released code can add a random term to each step (argument \texttt{eta}); it will be set to zero, so the teacher always takes the same path for the same input. The same velocity also gives, at no extra cost, the teacher's estimate of the final mel-spectrogram, found by moving the rest of the way in a single step:
\begin{equation}
\hat x_1(t_i) = x_{t_i} + (1-t_i)\,v(x_{t_i},t_i),
\label{eq:endpoint}
\end{equation}
\eqentry{Endpoint estimate at a teacher step}
\begin{whereterms}[2.6cm]
$\hat x_1(t_i)$ & endpoint estimate at grid point $i$: the teacher's current prediction of the final mel-spectrogram \\
$x_{t_i}$ & teacher state at time $t_i$ \\
$v(x_{t_i},t_i)$ & guided velocity at that state \\
$1-t_i$ & remaining distance to the speech end of the time scale
\end{whereterms}
This estimate is what the teacher hands to the student. The released system stops the teacher at the first grid point after its switching argument $s=0.25$ on the 16-interval grid, after eight steps, and hands over $\hat x_1(t_8)$. The student then mixes the handed-over estimate with fresh random noise, to the level $\tau$ at which it was trained to work, and predicts the output again. It repeats this once per point of its own schedule $\tau_j=1-\cos(\pi j/8)$, $j=0,\dots,3$, starting from position $j$, at one forward pass per step; the released code uses $j=1$, that is, three steps. Finally the prompt frames are removed and Vocos converts the mel-spectrogram into a waveform \parencite{siuzdak2024vocos}.

\begin{algorithm}[!htbp]
\caption{Released teacher-guided sampler of DMOSpeech~2 (baseline)}
\label{alg:baseline}
\small
\begin{algorithmic}[1]
\Require prompt mel and mask $M$, text, duration, stopping argument $s=0.25$, seed
\State draw $x_{t_0}\sim\mathcal{N}(0,I)$ with the input's seed
\State $k \gets \min\{i : t_i > s\}$ \Comment{16-interval sway grid}
\For{$i = 0,\dots,k-1$}
  \State $v \gets v_c(x_{t_i},t_i) + \gamma\,\big(v_c(x_{t_i},t_i) - v_u(x_{t_i},t_i)\big)$ \Comment{2 forward passes}
  \State $x_{t_{i+1}} \gets x_{t_i} + (t_{i+1}-t_i)\,v$ \Comment{Euler step}
\EndFor
\State $x \gets x_{t_k} + (1-t_k)\,v(x_{t_k},t_k)$ \Comment{endpoint estimate, 2 forward passes}
\For{$\tau \in \{\tau_1,\tau_2,\tau_3\}$} \Comment{student, $j=1$}
  \State draw $\varepsilon\sim\mathcal{N}(0,I)$ with a seed fixed per input and step
  \State $\phi \gets (1-\tau)\,\varepsilon + \tau\,x$; \quad $x \gets \phi + (1-\tau)\,g(\phi,\tau)$ \Comment{1 forward pass}
\EndFor
\State remove prompt frames; decode $x$ with the vocoder
\end{algorithmic}
\end{algorithm}

\phead{Candidate teacher paths} The router will run the teacher on the coarser 8-interval grid ($N=8$ in \Cref{eq:grid}), whose steps are twice as large as those of the released grid, and stop it after $k$ steps at the \emph{commit point} $t_k$, where it takes the estimate $\hat x_1=\hat x_1(t_k)$ of \Cref{eq:endpoint}. Three values of $k$ will be considered,
\begin{equation}
k\in\mathcal K=\{2,4,6\},\qquad t_k=1-\cos\Big(\frac{\pi k}{16}\Big),
\label{eq:paths}
\end{equation}
\eqentry{Candidate teacher paths}
\begin{whereterms}[2.6cm]
$k$ & number of guided Euler steps the teacher takes before the handover \\
$\mathcal K$ & set of candidate teacher paths, fixed before calibration \\
$t_k$ & commit point of path $k$ on the 8-interval grid
\end{whereterms}
These are the only points of the 8-interval grid that coincide with a point of the student's schedule, $\tau_{k/2}$, so the student can restart exactly where the teacher stops. The 8-interval grid consists of every second point of the 16-interval grid, so the four-step path reaches $t_4=t_8^{(16)}$, exactly the point where the released system hands over, in four steps of twice the size instead of eight; the two-step path hands over earlier and the six-step path later. Each path will be calibrated and evaluated separately (Sections~\ref{sec:obj3} and~\ref{sec:obj4}), and within one configuration the same path is used for every input. Before the payload is chosen, the student makes a quick first attempt at the whole output directly from noise, the draft,
\begin{equation}
s_0 = \varepsilon_0 + g(\varepsilon_0, 0),
\label{eq:draft}
\end{equation}
\eqentry{Student draft}
\begin{whereterms}[2.6cm]
$s_0$ & student draft: the student's one-step prediction of the output from noise \\
$\varepsilon_0$ & Gaussian noise drawn with the input's seed, $\varepsilon_0\sim\mathcal N(0,I)$ \\
$g(\varepsilon_0,0)$ & student network output at $\tau=0$ (one forward pass)
\end{whereterms}
which is the student step of Algorithm~\ref{alg:baseline} at $\tau=0$ and does not depend on the teacher. After routing, the student finishes from the selected payload starting at position $j=k/2$ of its schedule, that is, with $4-k/2$ steps from $\tau_{k/2}=t_k$. Its first step mixes in exactly the amount of noise that was still left when the teacher stopped, so the student continues from where the teacher left off. For $k=4$, the two-step finish is also the configuration of the DMOSpeech~2 paper \parencite{li2025dmospeech2}.

\phead{Cost} Counting forward passes gives the exact cost of every system in the study, which the rest of the chapter uses as the unit of computation.

\begin{proposition}[Cost identity]\label{prop:cost}
A system that makes $e$ guided teacher evaluations, $u\in\{0,1\}$ student draft passes and $n_s$ student steps costs
\begin{equation}
C = 2e + u + n_s
\label{eq:cost}
\end{equation}
\eqentry{Cost identity}
\begin{whereterms}[2.6cm]
$C$ & forward passes needed to synthesize one utterance \\
$e$ & guided teacher evaluations: Euler steps plus the endpoint evaluation \\
$2$ & forward passes per guided evaluation (conditional and unconditional) \\
$u$ & 1 if the student draft of \Cref{eq:draft} is computed, 0 otherwise \\
$n_s$ & student steps after the handover, $4-j$ for start index $j$
\end{whereterms}
forward passes. In particular, the released system costs $2\cdot9+0+3=21$ passes, and the router with teacher path $k$ costs
\[
C(k)=2(k+1)+1+\big(4-k/2\big)=\tfrac32k+7,
\]
for every input and whichever payload it selects; every candidate path costs less than the released system. The draft accounts for one of these passes even when the teacher's estimate is handed over, because the residual needs it.
\end{proposition}
\begin{proof}
Each guided evaluation runs the teacher once with and once without the condition; each student step and the draft run the student once (Algorithm~\ref{alg:baseline}, \Cref{eq:draft}). The released system makes $k=8$ Euler steps and one endpoint evaluation, $e=9$, with $j=1$. The router with path $k$ makes $k$ Euler steps and one endpoint evaluation, $e=k+1$, computes the draft, $u=1$, and finishes from $j=k/2$, so $n_s=4-k/2$. The quantile projection, the residual and the SNR estimate (Section~\ref{sec:obj2}) involve no network.
\end{proof}

Table~\ref{tab:cost} applies \Cref{eq:cost} to every system that will be compared. Cost is counted in forward passes because the count follows exactly from the procedure and does not depend on the computer used; measured time and real-time factor on the Tesla T4 GPU will be reported second, because they also include the duration predictor, the vocoder and the time spent moving data (Section~\ref{sec:measures}).

\begin{table}[!htbp]
\centering\small
\caption{Forward passes per utterance of the compared systems}
\label{tab:cost}
\begin{tabular}{@{}>{\raggedright\arraybackslash}p{4.8cm} c c c c@{}}
\toprule
\textbf{System} & \textbf{Teacher $2e$} & \textbf{Draft $u$} & \textbf{Student $n_s$} & \textbf{Total $C$} \\
\midrule
F5-TTS teacher alone (32 steps) & 64 & 0 & 0 & 64 \\
DMOSpeech~2, released system ($s=0.25$) & 18 & 0 & 3 & 21 \\
Payload router, path $k\in\{2,4,6\}$ & $2k+2$ & 1 & $4-k/2$ & $C(k)$ \\
Never project, path $k$ & $2k+2$ & 0 & $4-k/2$ & $C(k)-1$ \\
Always project, path $k$ & $2k+2$ & 1 & $4-k/2$ & $C(k)$ \\
SNR-only router, path $k$ & $2k+2$ & 0 or 1 & $4-k/2$ & $\le C(k)$ \\
Student restarted from its draft (\texttt{STUR2}) & 0 & 1 & 2 & 3 \\
DMOSpeech~2 student alone & 0 & 0 & 4 & 4 \\
\bottomrule
\end{tabular}

\end{table}

\subsubsection{Notation}
Table~\ref{tab:notation} lists the symbols used in the rest of this chapter. Inputs are treated as random draws from an unknown population $\mathcal{D}$ of (prompt, prompt text, target text, prompt condition) combinations; this is what allows the calibration to make statements about new inputs from the same population.

\begin{table}[!htbp]
\centering\small
\caption{Notation}
\label{tab:notation}
\begin{tabularx}{\columnwidth}{@{}>{\raggedright\arraybackslash}p{2.9cm} X@{}}
\toprule
Symbol & Meaning \\
\midrule
$x$ & input (prompt audio, prompt text, target text, prompt condition) drawn from distribution $\mathcal{D}$ \\
$k$, $\mathcal K$ & teacher path (number of teacher steps); candidate paths $\{2,4,6\}$ \\
$C(k)$ & forward passes per utterance of the router with path $k$ (Prop.~\ref{prop:cost}) \\
$\hat x_1$ & teacher endpoint estimate at the commit point $t_k$ (generated frames, $n\times B$) \\
$s_0$ & student draft (\Cref{eq:draft}) \\
$Q(\hat x_1,s_0)$ & quantile projection payload (Def.~\ref{def:quantile}) \\
$r$ & invariant residual (Def.~\ref{def:residual}) \\
$\hat S$, $\theta$ & estimated prompt SNR in dB; fixed SNR gate \\
$n_{\text{dev}}$, $n$, $n_{\text{test}}$ & numbers of development, calibration and test speakers \\
$\pi_{k,\lambda}$ & routing rule with path $k$ and residual threshold $\lambda$ (Def.~\ref{def:policy}) \\
$\Lambda=\{\lambda_1<\dots<\lambda_G\}$ & ordered grid of candidate thresholds \\
$\hat\lambda_k$ & calibrated threshold of path $k$ \\
$F(x,y)\in\{0,1\}$ & quality failure of output $y$ for input $x$ (Def.~\ref{def:failure}) \\
$y_{\text{ref}}(x)$, $y_{k,\lambda}(x)$ & output of the released system; output of the router with path $k$ and threshold $\lambda$ \\
$e(x,k,\lambda)$ & excess-failure indicator (\Cref{eq:excess}) \\
$R(k,\lambda)$ & risk: $\E_{x\sim\mathcal{D}}[e(x,k,\lambda)]$ \\
$\rho(k,\lambda)$ & routing rate: share of inputs that receive $Q$ \\
$\alpha,\delta$ & tolerated risk; allowed probability that calibration fails \\
$C$ & forward passes per utterance (Prop.~\ref{prop:cost}) \\
\bottomrule
\end{tabularx}
\end{table}

\subsubsection{Data and Splits}
\label{sec:data}

The corpus will be Seed-TTS test-en, whose prompts come from Common Voice \parencite{anastassiou2024seedtts,ardila2020commonvoice}; it is the English test set used by DMOSpeech~2 \parencite{li2025dmospeech2}, so the reproduced baseline can be compared with published numbers. It contains 1,088 utterances from 666 prompt speakers. The speakers will be divided at random into three groups, called splits, with no speaker in more than one split (Table~\ref{tab:splits}). The splits will use $n_{\text{dev}}+n+n_{\text{test}}$ of the 666 speakers and one sentence per speaker; speaker identity is taken from the prompt recording. The development split will serve the characterization stage and the fixing of the minimum quality levels; the calibration split will serve only the calibration of the threshold of each path; and the test split will serve the check of the guarantee and the comparisons. Keeping the speakers separate prevents the threshold from being tuned to voices it is later tested on, and one sentence per speaker keeps the inputs close to independent, which the guarantee of Proposition~\ref{prop:ltt} requires. The development runs of Appendix~\ref{app:snr-r} used development speakers only, so the calibration and test speakers will be new to the method.

A second test set, the test-clean portion of LibriSpeech-PC \parencite{meister2023librispeechpc}, will be used only for testing and never for development or calibration. Its audio is LibriSpeech test-clean \parencite{panayotov2015librispeech} and its transcripts are the LibriSpeech-PC versions with punctuation and capitalization restored. It differs from Seed-TTS test-en in source (read audiobooks rather than crowd-sourced recordings), so it checks whether the threshold calibrated on Seed-TTS speakers carries over to a new source. One utterance per speaker will be used, chosen in the manifest before testing. Because the guarantee covers only inputs like those used in calibration, results on this set will be reported as a check of how well the method transfers, not as part of the guarantee.

Before the splits are drawn, a list called the manifest will record every selected utterance with its speaker, transcripts, duration, split and assigned prompt condition, and it will not change after the experiments begin. An utterance will be excluded only for reasons fixed in advance, namely a missing transcript or audio that cannot be read, and every exclusion will be listed with its reason. A synthesis run that fails or produces unusable audio will be rerun once with the same seeds; if it fails again, it will be counted as a quality failure of that system rather than removed. If the released system itself fails to produce audio for an input, the input will be kept and its reference output will count as a failure.

\phead{Prompt conditions} Every prompt will be synthesized under eight conditions (Table~\ref{tab:conditions}). The clean condition uses the prompt as recorded. The other seven add a kind of damage that a real recording can have. The signal-to-noise ratio, in decibels (dB), states how much louder the speech is than the added noise; a lower value means a noisier prompt. White noise at 15~dB, a steady hiss clearly heard behind the speech, tests a clearly noisy prompt. Three moderate conditions at 25~dB, with white, pink and high-pass noise, add damage that is audible but that a simple loudness-based check can mistake for a clean recording, as the panel requested; the three noises differ in which frequencies carry their energy. Reverberation (echo of a room), a telephone line and background babble test damage that changes the prompt's timing or frequency range rather than adding a steady noise. All added noise is generated from a seed fixed for each input. Within an input, every system uses the same random starting values (Section~\ref{sec:baseline}), so differences between systems are paired.

\begin{table}[!htbp]
\centering\small
\caption{Prompt conditions}
\label{tab:conditions}
\begin{tabularx}{\textwidth}{@{}>{\raggedright\arraybackslash}p{2.6cm} >{\raggedright\arraybackslash}X >{\raggedright\arraybackslash}p{3.0cm}@{}}
\toprule
\textbf{Condition} & \textbf{How the prompt is changed} & \textbf{Kind} \\
\midrule
Clean & As recorded & None \\
White, 15~dB & White Gaussian noise at 15~dB & Additive, heavy \\
White, 25~dB & White Gaussian noise at 25~dB & Additive, moderate \\
Pink, 25~dB & Noise whose energy falls with frequency (a deeper rushing sound), at 25~dB & Additive, moderate \\
High-pass, 25~dB & Noise with its energy above about 2~kHz (hiss), at 25~dB & Additive, moderate \\
Reverberation & Simulated room echo that takes 0.9~s to die away, loudness matched & Echo \\
Telephone & Limited to the 300--3,400~Hz range of a telephone line, loudness matched & Narrow frequency range \\
Babble, 15~dB & Recorded cafeteria chatter from DEMAND \parencite{thiemann2013demand} at 15~dB & Additive, changing \\
\bottomrule
\end{tabularx}
\end{table}

All eight conditions will be synthesized for every development, calibration and test speaker, so that results can be reported per condition. For the calibration and for the check of the guarantee, however, each calibration and test speaker will contribute one input: the manifest will assign each speaker one condition at random, so that every condition goes to the same number of speakers, one eighth of each split. This keeps one observation per speaker, as the guarantee requires, and makes the calibration inputs an equal mix of the eight conditions. The guarantee therefore holds on average over that mix; per-condition rates will be reported separately (Section~\ref{sec:robust}).

\begin{table}[!htbp]
\centering\small
\caption{Data splits, with no speaker in more than one split}
\label{tab:splits}
\begin{tabularx}{\textwidth}{@{}>{\raggedright\arraybackslash}p{2.6cm} >{\raggedright\arraybackslash}X >{\raggedright\arraybackslash}p{3.6cm} >{\raggedright\arraybackslash}p{2.5cm}@{}}
\toprule
\textbf{Split} & \textbf{Used for} & \textbf{Constraint} & \textbf{Size} \\
\midrule
Development & Characterization; minimum quality levels $s_{\min}$, $u_{\min}$ & No speaker shared with other splits & $n_{\text{dev}}$ speakers, one utterance each \\
Calibration & Threshold $\hat\lambda_k$ of each path & One utterance per speaker & $n$ speakers, one utterance each \\
Test & Check of the guarantee; comparisons & One utterance per speaker & $n_{\text{test}}$ speakers, one utterance each \\
LibriSpeech-PC test-clean & Transfer to a new source & Test only; one utterance per speaker & $n_{\text{LS}}$ speakers, one utterance each \\
\bottomrule
\end{tabularx}
\end{table}

\subsubsection{Quality, Variety, and Cost Measures}
\label{sec:measures}

Each generated sentence will be scored on three qualities (Table~\ref{tab:measures}): whether its words can be understood, whether it sounds like the intended speaker, and how natural it sounds. Intelligibility is measured by the word error rate (WER): a speech recognizer, Whisper-large-v3, writes down what it hears, and the transcript is compared with the intended text, as in the DMOSpeech~2 evaluation \parencite{radford2023whisper,li2025dmospeech2}. Before comparison, both texts are put in the same form as in the Seed-TTS English evaluation: punctuation is replaced by spaces, repeated spaces are removed and all letters are lowercased. Speaker similarity (SIM) is measured with a speaker-verification model, WavLM-Base+ (\texttt{microsoft/wavlm-base-plus-sv}) \parencite{chen2022wavlm}, which turns a recording into a list of numbers, an embedding, that describes the voice; SIM is the cosine similarity of the embeddings of the generated speech and of the real recording of the target sentence, where 1 means the same voice. DMOSpeech~2 uses a larger verifier, so the SIM values of this study are not directly comparable with its published numbers; every comparison within the study uses the same verifier. Naturalness will be predicted with UTMOS \parencite{saeki2022utmos}, a network trained to predict the score from 1 to 5 that human listeners would give. Word error rate counts the word-level edits needed to turn the recognized transcript into the target text:
\begin{equation}
\mathrm{WER} = \frac{S_w + D_w + I_w}{N_w}.
\label{eq:wer}
\end{equation}
\eqentry{Word error rate}
\begin{whereterms}[2.6cm]
$S_w, D_w, I_w$ & numbers of substituted, deleted and inserted words in the transcript \\
$N_w$ & number of words in the normalized target text
\end{whereterms}
Because inserted words are counted, WER can exceed one when the output keeps talking past the end of the text; such run-on outputs will be counted and reported separately. Speaker similarity compares the speaker embeddings of the generated and the ground-truth speech:
\begin{equation}
\mathrm{SIM} = \frac{e_{\text{gen}}^{\top} e_{\text{tgt}}}{\lVert e_{\text{gen}}\rVert\,\lVert e_{\text{tgt}}\rVert}.
\label{eq:sim}
\end{equation}
\eqentry{Speaker similarity}
\begin{whereterms}[2.6cm]
$e_{\text{gen}}$ & speaker embedding of the generated speech \\
$e_{\text{tgt}}$ & speaker embedding of the ground-truth target recording \\
$\lVert\cdot\rVert$ & Euclidean length of a vector
\end{whereterms}

\phead{Prosodic variety} A text-to-speech system that gives the same melody every time it reads a sentence sounds monotonous, and the teacher-guided generation of DMOSpeech~2 is meant in part to recover variety that the distilled student loses \parencite{li2025dmospeech2}. Variety will therefore be measured, because it is the reason the teacher is kept at all (Section~\ref{sec:obj2}). Each input is synthesized $K$ times with different random seeds, and the pitch (fundamental frequency) of each output is tracked frame by frame with pYIN \parencite{mauch2014pyin}. At each frame, the spread of the pitch across the seeds is divided by its average (the coefficient of variation), and the result is averaged over the frames in which enough of the outputs are voiced:
\begin{equation}
\mathrm{CV}_{f_0} = \frac{1}{\lvert\mathcal V\rvert}\sum_{t\in\mathcal V}\frac{\operatorname{sd}_k\big(f_0^{(k)}(t)\big)}{\operatorname{mean}_k\big(f_0^{(k)}(t)\big)}.
\label{eq:cvf0}
\end{equation}
\eqentry{Pitch variety across seeds}
\begin{whereterms}[2.6cm]
$\mathrm{CV}_{f_0}$ & pitch variety of one system on one input; higher means more varied melodies \\
$f_0^{(k)}(t)$ & fundamental frequency of the output of seed $k$ at frame $t$ \\
$K$ & number of seeds \\
$\operatorname{sd}_k$, $\operatorname{mean}_k$ & standard deviation and mean across the seeds that voice frame $t$ \\
$\mathcal V$ & frames voiced by at least $\max(2,\lceil0.3K\rceil)$ seeds
\end{whereterms}
Because different seeds also change the timing of the words, an aligned version will be reported as well, in which every output is first stretched in time to line up with the first one by dynamic time warping \parencite{sakoe1978dtw}; it separates variety in melody from variety in timing. Pitch-based measures capture only part of what listeners hear as varied prosody, since they ignore rhythm and loudness \parencite{yang2025prosody}, and newer neural pitch trackers are more accurate than pYIN \parencite{morrison2023pitch}; $\mathrm{CV}_{f_0}$ is kept because it is the measure of the team's development runs, and its limits will be stated with the results. Variety will be measured on $n_{\text{var}}$ test speakers under the clean condition and one moderate noise condition, with the same $K$ seeds for every system.

\phead{Cost} Cost will be measured primarily in forward passes (Proposition~\ref{prop:cost}) and secondarily in measured time and real-time factor on the study's GPU, an NVIDIA Tesla T4 through Google Colab. The real-time factor relates synthesis time to the length of the speech produced; values below one mean faster than real time:
\begin{equation}
\mathrm{RTF} = \frac{T_{\text{synth}}}{T_{\text{audio}}}.
\label{eq:rtf}
\end{equation}
\eqentry{Real-time factor}
\begin{whereterms}[2.6cm]
$T_{\text{synth}}$ & wall-clock time of the full synthesis pipeline for one utterance \\
$T_{\text{audio}}$ & duration of the generated speech
\end{whereterms}

\begin{table}[!htbp]
\centering\small
\caption{Measures of quality, variety, risk and cost}
\label{tab:measures}
\label{tab:matrix}
\label{tab:matrix-cols}
\begin{tabularx}{\textwidth}{@{}>{\raggedright\arraybackslash}p{2.4cm} >{\raggedright\arraybackslash}X >{\raggedright\arraybackslash}p{3.0cm} >{\raggedright\arraybackslash}p{3.0cm}@{}}
\toprule
\textbf{Measure} & \textbf{Computed with} & \textbf{Better when} & \textbf{Role} \\
\midrule
WER & Whisper-large-v3 transcript, stated normalization (\Cref{eq:wer}) & Lower & Failure rule; compared \\
SIM & WavLM-Base+ speaker verifier, cosine similarity (\Cref{eq:sim}) & Higher & Failure rule; compared \\
UTMOS & Predicted mean opinion score & Higher & Failure rule; compared \\
Prosodic variety, $\mathrm{CV}_{f_0}$ & Spread of pitch across $K$ seeds of the same sentence, raw and time-aligned (\Cref{eq:cvf0}) & Higher & Compared; main test of the teacher's value \\
Failure rate & Share of inputs that fail the quality rule, with an exact interval (Def.~\ref{def:failure}) & Lower & Compared \\
Excess-failure rate & Share of inputs that fail while the released system passes, with an exact interval (\Cref{eq:excess}) & Lower & Calibration; compared \\
Routing rate & Share of inputs given the projection & --- & Router only \\
Forward passes, saving & Proposition~\ref{prop:cost}; share of the released system's passes saved (\Cref{eq:saving}) & Fewer passes & Primary cost \\
Runtime, RTF & Wall-clock time on the Tesla T4 GPU, full pipeline (\Cref{eq:rtf}) & Lower & Secondary cost \\
\bottomrule
\end{tabularx}
\end{table}

\phead{Quality failure} The calibration needs a yes-or-no outcome, failed or not, that can be counted for every system on every input. The panel asked that the baseline not be treated as failure-free by definition, so the rule will judge every output by the same fixed standard: the same rule will be applied to the released system, to the router and to every other system.

\begin{definition}[Quality failure]\label{def:failure}
An output $y$ for input $x$ fails, $F(x,y)=1$, if
\begin{equation}
\mathrm{WER}(x,y) > c
\quad\text{or}\quad
\mathrm{SIM}(x,y) < s_{\min}
\quad\text{or}\quad
\mathrm{UTMOS}(y) < u_{\min},
\label{eq:failure}
\end{equation}
\eqentry{Quality-failure rule}
and $F(x,y)=0$ otherwise.
\begin{whereterms}[2.6cm]
$F(x,y)$ & quality-failure indicator of output $y$ on input $x$ \\
$\mathrm{WER}(x,y)$, $\mathrm{SIM}(x,y)$, $\mathrm{UTMOS}(y)$ & the three measures of output $y$ \\
$c$ & WER ceiling \\
$s_{\min}$ & SIM floor (\Cref{eq:floors}) \\
$u_{\min}$ & UTMOS floor (\Cref{eq:floors})
\end{whereterms}
\end{definition}

The WER ceiling $c$ is set so that a short sentence does not fail because of a single recognizer error. Minimum SIM and UTMOS levels cannot be taken from the literature, because both scores depend on the scoring model and the prompts. They will therefore be set from the reproduced baseline on the development speakers, before any calibration:
\begin{equation}
\begin{aligned}
s_{\min} &= q_{p}\big\{\mathrm{SIM}(x, y_{\text{ref}}(x)) : x\in\mathcal X_{\text{dev}}^{\text{clean}}\big\},\\
u_{\min} &= q_{p}\big\{\mathrm{UTMOS}(y_{\text{ref}}(x)) : x\in\mathcal X_{\text{dev}}^{\text{clean}}\big\}.
\end{aligned}
\label{eq:floors}
\end{equation}
\eqentry{SIM and UTMOS floors}
\begin{whereterms}[2.6cm]
$q_{p}\{\cdot\}$ & $p$-quantile of a set of scores: the value below which a share $p$ of the scores lies \\
$y_{\text{ref}}(x)$ & output of the released DMOSpeech~2 system for input $x$ \\
$\mathcal X_{\text{dev}}^{\text{clean}}$ & development inputs with clean prompts
\end{whereterms}
In plain words, an output fails on similarity or naturalness when it is worse than all but the worst share $p$ of what the released system produces from clean prompts. The floors describe what the baseline already achieves rather than an ideal, so the released system will itself fail on some inputs, and its failure rate will be measured and reported like every other system's. UTMOS is included in the rule, unlike in the earlier proposal, because the router's choice mainly affects naturalness on damaged prompts (Appendix~\ref{app:snr-r}), and a rule without it would not see the failures the router is meant to prevent. Any output above the WER ceiling will also be checked by its transcript and by listening before the failure is attributed to synthesis rather than to the recognizer. All constants of the rule will be fixed before calibration and reported with the results.

A failed output is not removed from the data. The failure rate of a system is the number of inputs on which it fails divided by the number of inputs, reported with an exact Clopper--Pearson interval \parencite{clopper1934}, the range within which the true rate lies with 95\% confidence.

\subsubsection{Tools, Hardware, and Reproducibility}
\label{sec:tools}

The implementation will use Python and PyTorch with the released DMOSpeech~2 code (Section~\ref{sec:baseline}), Vocos \parencite{siuzdak2024vocos} and the evaluation models above on an NVIDIA Tesla T4 GPU (15.6~GB of memory) through Google Colab, in 32-bit floating point with seeded random generators. Timing will cover the whole generation process (prompt SNR estimate, teacher steps, student draft, routing, student steps and vocoder) and will be measured for each candidate path. The clock will be read only after the GPU has finished its work. Five sentences will be generated before timing starts, so that start-up costs are not counted, and the order of the compared systems will be varied across sentences. Preprocessing, duration prediction, model loading, file writing and the offline quality evaluation will be reported separately. The GPU model and software versions of every Colab session will be recorded. Every script will use one shared generation module, so all systems run the same code. Before any experiment, that module will be checked for identical outputs from identical seeds, for the defining properties of the projection and the residual, for exact forward-pass counts (Table~\ref{tab:cost}), and for an exact reproduction of the released system by the fallback. Raw per-utterance results and configuration files will be archived, and every table will be regenerated by a committed script from those raw results.

% =====================================================================
\subsection{Characterizing the Baseline and the Handover Payloads}
\label{sec:obj1}

This stage covers Steps~4 to~6 of Table~\ref{tab:steps} and uses the development speakers only. Its purpose is to measure, before the method is calibrated, what the established systems achieve on the study's prompts and what each of the two payloads does to quality.

\phead{Reproducing the baseline} The released DMOSpeech~2 system will be run on every development input under all eight prompt conditions, together with the F5-TTS teacher alone and the DMOSpeech~2 student alone (Table~\ref{tab:cost}). Their WER, SIM and UTMOS will be reported as absolute values per condition, with the number of outputs that exceed the WER ceiling and the number of run-on outputs. As a reproduction check, the released system will also be run once on the full Seed-TTS test-en set with clean prompts, and its WER will be compared with the published value \parencite{li2025dmospeech2}; SIM will be reported with the caveat on verifiers of Section~\ref{sec:measures}. Nothing is tuned on this run: no threshold or minimum level is set from it, so the calibration and test speakers it includes do not influence the design.

\phead{Characterizing the payloads} On the same inputs, every candidate teacher path will be run, and both of its payloads will be completed by the student: the teacher payload $\hat x_1$ (never project) and the quantile projection $Q(\hat x_1,s_0)$ (always project). One run of the six-step path is enough, because it passes through $t_2$ and $t_4$ and the velocity it computes at those points gives their endpoint estimates (\Cref{eq:endpoint}) without extra passes. The student restarted from its own draft (\texttt{STUR2}) will be run as well, because it uses the same draft and the same two-step finish without any teacher. For every input and path, the prompt SNR $\hat S$ and the residual $r$ will be recorded. The analysis will report, per condition and path, the cost and quality of each payload, how often each payload fails under Definition~\ref{def:failure}, and how $r$ and $\hat S$ relate to the cases in which the teacher payload fails and the projection does not. The numerical identities of Section~\ref{sec:tools} will also be checked on real outputs, including that the projection stays the same when a steady noise or a change of loudness is applied to the teacher's estimate (Proposition~\ref{prop:invariance}).

\phead{Outputs} This stage fixes the minimum SIM and UTMOS levels (\Cref{eq:floors}) and produces the reproduced reference values that every later comparison uses. It does not choose the teacher path or the residual threshold, which are set only by the calibration of Section~\ref{sec:obj3}.

Algorithm~\ref{alg:grid} gives the procedure that records, for one input, the outputs and scores that the characterization, the calibration and the evaluation all use. Because all payloads come from one teacher run and one draft, and because the routing decision depends only on recorded values, the router's output for any path and threshold can be read from this record without generating the speech again.

\begin{algorithm}[!htbp]
\caption{Recording the outcomes of one input}
\label{alg:grid}
\small
\begin{algorithmic}[1]
\Require input $x$ with its prompt condition and seed
\State run the released system (Algorithm~\ref{alg:baseline}); store $y_{\text{ref}}$ and $F(x,y_{\text{ref}})$
\State run the teacher on the 8-interval grid to $t_6$; store $\hat x_1^{(k)}=\hat x_1(t_k)$ for $k\in\mathcal K$
\State compute the draft $s_0$ (\Cref{eq:draft}) and $\hat S$ (\Cref{eq:snr}) \Comment{1 pass}
\For{$k \in \mathcal K$}
  \State compute $Q(\hat x_1^{(k)},s_0)$ (\Cref{eq:quantile}) and $r^{(k)}$ (\Cref{eq:residual}) \Comment{no pass}
  \For{$P \in \{\hat x_1^{(k)},\ Q(\hat x_1^{(k)},s_0)\}$}
    \State finish with $4-k/2$ student steps from $P$ ($j=k/2$, shared noise draws)
    \State decode, score WER, SIM and UTMOS; store the output and $F$
  \EndFor
\EndFor
\State store $\hat S$, every $r^{(k)}$ and all scores in the input's record
\end{algorithmic}
\end{algorithm}

% =====================================================================
\subsection{Formulation and Analysis of the Payload Router}
\label{sec:obj2}

This section covers Step~7 of Table~\ref{tab:steps}. It states the problem the router solves, defines its parts, and proves the properties on which the design rests. It uses no data.

\subsubsection{Problem Formulation}
\label{sec:problem-formulation}

Once a teacher path $k$ is fixed, every input costs $C(k)$ passes (Proposition~\ref{prop:cost}), and the router only chooses the payload. The teacher payload carries all of the teacher's detail, including its variety of pitch and timing; the projection keeps only the teacher's timing and takes its loudness values from the student draft (Section~\ref{sec:quantile}), so it is safer when the prompt is damaged but carries less of the teacher. The router should therefore pass the teacher payload whenever it is safe to do so. The threshold $\lambda$ of the routing rule $\pi_{k,\lambda}$ (Definition~\ref{def:policy}) controls this: the higher the threshold, the more inputs receive the teacher payload. For each path, the problem is to find the highest threshold that keeps the risk within the tolerance $\alpha$:
\begin{equation}
\lambda^\ast_k=\max\big\{\lambda\in\Lambda : R(k,\lambda)\le\alpha\big\},\qquad k\in\mathcal K,
\label{eq:objective}
\end{equation}
\eqentry{Router calibration problem}
\begin{whereterms}[2.6cm]
$\lambda^\ast_k$ & highest safe threshold for path $k$; a larger $\lambda$ projects fewer inputs \\
$\mathcal K$, $\Lambda$ & candidate paths (\Cref{eq:paths}) and finite, ordered grid of candidate thresholds (Section~\ref{sec:calibration}) \\
$R(k,\lambda)$ & risk of the router with path $k$ and threshold $\lambda$: $\E_{x\sim\mathcal D}\big[e(x,k,\lambda)\big]$ \\
$e(x,k,\lambda)$ & excess-failure indicator (\Cref{eq:excess}) \\
$\E_{x\sim\mathcal D}$ & average over new inputs drawn from the population $\mathcal D$ \\
$\alpha$ & tolerated risk
\end{whereterms}
The risks are unknown, so \Cref{eq:objective} cannot be solved exactly; the calibration of Section~\ref{sec:obj3} estimates $\lambda^\ast_k$ from data, with a stated confidence. If no threshold is safe for a path, that path falls back to the released system. The paths are not compared here: each is calibrated on its own, and the comparison of cost, quality and risk across paths is part of the evaluation (Section~\ref{sec:obj4}).

The risk is the share of inputs on which the router fails but the released system does not. For one input this event, an \emph{excess failure}, is written
\begin{equation}
e(x,k,\lambda) = F\big(x,y_{k,\lambda}(x)\big)\,\big(1-F\big(x,y_{\text{ref}}(x)\big)\big).
\label{eq:excess}
\end{equation}
\eqentry{Excess-failure indicator}
\begin{whereterms}[2.6cm]
$e(x,k,\lambda)$ & 1 if the router with path $k$ and threshold $\lambda$ fails on $x$ while the released system passes, 0 otherwise \\
$y_{k,\lambda}(x)$ & router output for $x$ with path $k$ and threshold $\lambda$ \\
$y_{\text{ref}}(x)$ & released DMOSpeech~2 output for $x$ \\
$F$ & quality-failure indicator (Definition~\ref{def:failure})
\end{whereterms}
Because the routing rate $\rho(k,\lambda)$, the share of inputs that receive the projection, can only fall as $\lambda$ rises (Proposition~\ref{prop:mono}), choosing the highest safe $\lambda$ keeps as many teacher payloads as the risk allows on each path. Limiting the excess failures also limits how much more often the router fails than the released system:

\begin{proposition}[Failure rate relative to the baseline]\label{prop:rate}
Write $F_{k,\lambda}=F(x,y_{k,\lambda}(x))$ and $F_{\text{ref}}=F(x,y_{\text{ref}}(x))$. For every path $k$ and threshold $\lambda$,
\begin{equation}
\P(F_{k,\lambda}=1) \ \le\ \P(F_{\text{ref}}=1) + R(k,\lambda).
\label{eq:rate}
\end{equation}
\eqentry{Failure rate relative to the baseline}
\begin{whereterms}[2.6cm]
$\P(F_{k,\lambda}=1)$ & failure rate of the router on new inputs \\
$\P(F_{\text{ref}}=1)$ & failure rate of the released system on the same inputs \\
$R(k,\lambda)$ & risk of \Cref{eq:objective}
\end{whereterms}
Hence $R(k,\lambda)\le\alpha$ implies that the router fails on at most $\alpha$ more of the inputs than the released system does.
\end{proposition}
\begin{proof}
For every input, $F_{k,\lambda} \le F_{\text{ref}} + F_{k,\lambda}(1-F_{\text{ref}})$: if $F_{\text{ref}}=1$ the right side is at least 1, and if $F_{\text{ref}}=0$ both sides equal $F_{k,\lambda}$. Taking expectations over $x\sim\mathcal D$ gives \Cref{eq:rate}.
\end{proof}

In other words, if the router adds a failure on at most a share $\alpha$ of inputs, its overall failure rate can exceed that of the released system by at most $\alpha$. The comparison is one-sided on purpose: an input on which the router passes and the released system fails does not lower the risk, so the calibration cannot balance new failures on some inputs against improvements on others.

\subsubsection{Condition Model}
\label{sec:condition-model}

The projection rests on a simple model of how a damaged prompt reaches the teacher's output. Let $\hat x_1\in\R^{n\times B}$ hold the values of the $n$ generated frames in $B=100$ frequency bands, on a logarithmic loudness scale, and write $\hat x_{1,b}\in\R^n$ for band $b$. The earlier experiments found that the teacher copies the prompt's recording conditions, such as its background noise, into its estimate (Appendices~\ref{app:prelim} and~\ref{app:snr-r}). A recording channel that makes a band louder or softer by a fixed factor $h_b>0$, together with a steady background noise of power $\nu_b\ge0$, changes a value $z$ of that band approximately as
\begin{equation}
g_b(z) = \tfrac12\log\!\big(h_b\,e^{2z}+\nu_b\big),\qquad b=1,\dots,B.
\label{eq:condition}
\end{equation}
\eqentry{Stationary condition model}
\begin{whereterms}[2.6cm]
$g_b$ & effect of the recording condition on a log-mel value in band $b$ \\
$z$ & log-mel value of a frame in band $b$ (log of the amplitude) \\
$h_b$ & stationary gain of the channel in band $b$ \\
$\nu_b$ & power of a stationary noise floor in band $b$
\end{whereterms}
Every such $g_b$ is strictly increasing: a louder frame stays louder after the change. Changes of this kind can be arranged into three groups, each containing the one before, $G_1\subset G_2\subset G_3$: adding a constant to each band, $z\mapsto z+c_b$ (a change of volume alone); stretching and shifting each band, $z\mapsto a_bz+c_b$ with $a_b>0$; and any strictly increasing change of each band, which includes \Cref{eq:condition} for every channel and noise level. A larger group covers more recording conditions, but less of the teacher's estimate survives all of them unchanged. For $G_3$, what survives is the time order within each band: such a change alters the values but never which frame of a band is louder than which. Table~\ref{tab:ladder} lists, for each group, the part of the estimate that the group cannot change and the payload built from it.

\begin{table}[!htbp]
\centering\small
\caption{Condition groups and the payloads that are unchanged by them}
\label{tab:ladder}
\begin{tabularx}{\textwidth}{@{}>{\raggedright\arraybackslash}p{3.2cm} >{\raggedright\arraybackslash}p{3.5cm} >{\raggedright\arraybackslash}X@{}}
\toprule
\textbf{Group} & \textbf{Unchanged part of $\hat x_{1,b}$} & \textbf{Payload} \\
\midrule
None & The whole estimate & Teacher payload $\hat x_1$ \\
$G_1$: per-band shifts & Values minus the band mean & Teacher's centred band plus the draft's band mean \\
$G_2$: per-band affine maps & Standardized values ($z$-scores) & Teacher's $z$-scores scaled to the draft's band mean and spread \\
$G_3$: per-band increasing maps & Time order (ranks) & Quantile projection $Q(\hat x_1,s_0)$ \\
\bottomrule
\end{tabularx}
\end{table}

The model has limits that will be kept in view. Real background noise varies from frame to frame, so it is not exactly a fixed change; and echo, telephone filtering and babble change the prompt's timing or frequency content in ways that no change of single bands describes. The model therefore explains why the projection should remove steady background noise and channel effects, and the eight prompt conditions test how far it holds beyond them.

\subsubsection{Quantile Projection}
\label{sec:quantile}

\begin{definition}[Quantile projection]\label{def:quantile}
For each band $b$, let $\operatorname{rank}_b(t)\in\{1,\dots,n\}$ be the position of the teacher's value $\hat x_{1,b}(t)$ when the $n$ generated frames of band $b$ are sorted from quietest to loudest, with ties broken by frame order, and let $s_{0,b,(1)}\le\dots\le s_{0,b,(n)}$ be the sorted values of the draft in the same band. The quantile projection is
\begin{equation}
Q(\hat x_1,s_0)_{t,b} = s_{0,b,(\operatorname{rank}_b(t))},\qquad t=1,\dots,n,\ b=1,\dots,B,
\label{eq:quantile}
\end{equation}
\eqentry{Quantile projection}
on the generated frames; the prompt frames of the payload are taken from the draft.
\begin{whereterms}[2.6cm]
$Q(\hat x_1,s_0)_{t,b}$ & value of the projected payload at frame $t$ in band $b$ \\
$\operatorname{rank}_b(t)$ & rank of the teacher's value at frame $t$ within band $b$ \\
$s_{0,b,(k)}$ & $k$-th smallest value of the draft in band $b$ \\
$n$, $B$ & numbers of generated frames and of mel bands
\end{whereterms}
\end{definition}

In words, in every band the projection puts the draft's values in the teacher's time order: the frame where the teacher is loudest receives the draft's largest value, the frame where it is second loudest receives the draft's second largest value, and so on. The draft supplies \emph{which} loudness values the band contains, and the student, which works from noise, does not copy the prompt's background noise in the same way; the teacher supplies \emph{when} they occur, which carries its timing and intonation. The operation is histogram matching, introduced for image enhancement \parencite{hall1974histogram} and used since 2000 to remove the mismatch that noise causes in the features of speech recognizers \parencite{dharanipragada2000nonlinear,delatorre2005histogram,hilger2006quantile}. What is new here is its use as the payload that a teacher hands to a student in text-to-speech, together with the residual it defines and the calibrated choice of when to apply it. It costs one sort per band and no network.

\begin{proposition}[Invariance and maximality of the projection]\label{prop:invariance}
Assume the values within each band are distinct. A payload $H(\hat x_1,s_0)$ satisfies $H(g(\hat x_1),s_0)=H(\hat x_1,s_0)$ for every $g\in G_3$ if and only if it depends on $\hat x_1$ only through its per-band ranks. Consequently $Q$ is unchanged by every condition in $G_3$; among such payloads it keeps all the information about $\hat x_1$ that is not changed by them; every value it contains is a value of the draft in the same band; and $Q(\hat x_1,s_0)=s_0$ whenever $\hat x_{1,b}=\varphi_b(s_{0,b})$ for strictly increasing maps $\varphi_b$.
\end{proposition}
\begin{proof}
A strictly increasing $g_b$ preserves every comparison $\hat x_{1,b}(t)<\hat x_{1,b}(t')$, so the ranks of $g_b(\hat x_{1,b})$ equal those of $\hat x_{1,b}$, and \Cref{eq:quantile} gives $Q(g(\hat x_1),s_0)=Q(\hat x_1,s_0)$. Conversely, two bands with the same ranks are linked by a strictly increasing map (the piecewise-linear map through their sorted value pairs), so a payload that is unchanged by $G_3$ must take the same value on both; it can depend on $\hat x_1$ only through the ranks. Because the draft's values in a band are distinct, different rank vectors give different outputs of \Cref{eq:quantile}, so $Q$ loses none of the rank information. Every entry of \Cref{eq:quantile} is some $s_{0,b,(k)}$. Finally, if $\hat x_{1,b}=\varphi_b(s_{0,b})$ with $\varphi_b$ increasing, the teacher's ranks equal the draft's ranks, and placing the draft's sorted values at its own ranks returns $s_{0,b}$.
\end{proof}

Put simply, the projection is unaffected by any steady noise or channel effect that the teacher copied from the prompt, and no other payload with this property keeps more of the teacher's estimate. This explains both the use and the cost of the projection. It removes the copied noise, because such conditions cannot change the time order; but it also removes everything about the teacher's estimate except its time order, including any loudness detail that made the teacher's output more expressive. This loss is why the router will not project every input (\Cref{eq:objective}).

\subsubsection{Invariant Residual}
\label{sec:residual}

The router needs to know, for each input, whether the teacher's loudness values have moved away from the draft's, as they do when the teacher copies background noise. The projection makes this comparison easy: it is the teacher's estimate with the draft's values put in place, so the difference between the two measures exactly the change in values.

\begin{definition}[Invariant residual]\label{def:residual}
\begin{equation}
r(\hat x_1,s_0) = \frac{\lVert \hat x_1 - Q(\hat x_1,s_0)\rVert_M}{\lVert \hat x_1\rVert_M + \epsilon}.
\label{eq:residual}
\end{equation}
\eqentry{Invariant residual}
\begin{whereterms}[2.6cm]
$r(\hat x_1,s_0)$ & invariant residual: relative distance between the teacher estimate and its projection \\
$\lVert\cdot\rVert_M$ & overall size of a mel-spectrogram: the square root of the sum of its squared values over the generated frames and all bands \\
$\epsilon$ & small constant against division by zero, $\epsilon=10^{-12}$
\end{whereterms}
\end{definition}

\begin{proposition}[Zero residual]\label{prop:null}
$r(\hat x_1,s_0)=0$ if and only if, in every band, the teacher's values are a reordering in time of the draft's values. In particular, differences in timing alone never raise the residual.
\end{proposition}
\begin{proof}
The denominator of \Cref{eq:residual} is positive, so $r=0$ exactly when $\hat x_1=Q(\hat x_1,s_0)$ on the generated frames. Each band of $Q$ is a reordering of the draft's values, so equality requires the two bands to hold the same values. Conversely, if they do, assigning the $k$-th smallest draft value to the frame that holds the $k$-th smallest teacher value reproduces the teacher's band, so the numerator is zero.
\end{proof}

\begin{proposition}[Residual as a distance between value distributions]\label{prop:w2}
The squared numerator of \Cref{eq:residual} equals
\begin{equation}
\lVert \hat x_1 - Q(\hat x_1,s_0)\rVert_M^2
= \sum_{b=1}^{B}\sum_{k=1}^{n}\big(\hat x_{1,b,(k)} - s_{0,b,(k)}\big)^2
= n\sum_{b=1}^{B} W_2^2\big(\hat\mu_{b},\,\hat\nu_{b}\big).
\label{eq:w2}
\end{equation}
\eqentry{Residual and the Wasserstein distance}
\begin{whereterms}[2.6cm]
$\hat x_{1,b,(k)}$, $s_{0,b,(k)}$ & $k$-th smallest value of the teacher estimate and of the draft in band $b$ \\
$\hat\mu_b$, $\hat\nu_b$ & empirical distributions of the values in band $b$ of the teacher estimate and of the draft \\
$W_2$ & 2-Wasserstein distance between two distributions on the real line
\end{whereterms}
Moreover, if the draft's sorted values lie at or below the teacher's in a band, adding steady background noise of power $\nu_b>0$ to the teacher estimate leaves $Q$ unchanged and strictly increases the numerator.
\end{proposition}
\begin{proof}
In band $b$, $Q$ places $s_{0,b,(k)}$ at the frame that holds $\hat x_{1,b,(k)}$, so the squared difference summed over frames is $\sum_k(\hat x_{1,b,(k)}-s_{0,b,(k)})^2$. For two empirical distributions with $n$ equally weighted points on the real line, the optimal coupling pairs order statistics \parencite{peyre2019computational}, so this sum is $nW_2^2(\hat\mu_b,\hat\nu_b)$. A noise floor acts through \Cref{eq:condition} with $h_b=1$, which is strictly increasing, so $Q$ is unchanged by Proposition~\ref{prop:invariance}; and $g_b(z)=\tfrac12\log(e^{2z}+\nu_b)>z$, so every teacher order statistic moves up, and every term $(g_b(\hat x_{1,b,(k)})-s_{0,b,(k)})^2$ with $\hat x_{1,b,(k)}\ge s_{0,b,(k)}$ grows.
\end{proof}

The Wasserstein distance of \Cref{eq:w2} measures how far apart two sets of values are: in each band, it compares the quietest teacher value with the quietest draft value, the second quietest with the second quietest, and so on. Together, the three propositions describe what the residual detects. It ignores timing (Proposition~\ref{prop:null}); it measures how far the teacher's range of loudness values in each band has moved from the draft's (\Cref{eq:w2}); and steady background noise copied from the prompt pushes it upward. Unlike the prompt SNR estimate, it looks at what the teacher actually produced rather than at the prompt, so it can respond to damage that the prompt check mistakes for a clean recording.

\subsubsection{Prompt SNR Estimate}
\label{sec:snr}

\begin{definition}[Prompt SNR estimate]\label{def:snr}
The prompt recording is cut into overlapping pieces 25~ms long, one every 10~ms, and the energy of each piece is measured. The loudest pieces are taken as speech and the quietest as background:
\begin{equation}
\hat S = \mathrm{clip}_{[0,50]}\!\left(10\log_{10}\frac{\bar E_{\text{top}}}{\bar E_{\text{bottom}}}\right).
\label{eq:snr}
\end{equation}
\eqentry{Prompt SNR estimate}
\begin{whereterms}[2.6cm]
$\hat S$ & estimated prompt signal-to-noise ratio in dB \\
$\bar E_{\text{top}}$ & mean energy of the loudest half of the frames (speech) \\
$\bar E_{\text{bottom}}$ & mean energy of the quietest 10\% of the frames (background) \\
$\mathrm{clip}_{[0,50]}$ & limits the value to between 0 and 50~dB
\end{whereterms}
It is computed on the prompt before its volume is adjusted, and set to 0~dB when the prompt is shorter than two pieces.
\end{definition}

The gate $\theta$ will be fixed before calibration, at the value used in the earlier experiments (Appendix~\ref{app:prelim}), rather than calibrated. Keeping it fixed leaves the calibration with only the path and the residual threshold to choose, so the guarantee of Section~\ref{sec:obj3} applies to exactly the rule that will be used. The gate serves as a cheap first check for prompts that are clearly noisy; prompts that it passes are left to the residual.

\subsubsection{Routing Rule and Properties}
\label{sec:properties}

\begin{definition}[Payload routing rule]\label{def:policy}
For path $k$ and threshold $\lambda$, with $\hat x_1$ the teacher's estimate at $t_k$, the router hands the student
\begin{equation}
\pi_{k,\lambda}(x) =
\begin{cases}
Q(\hat x_1,s_0) & \text{if } \hat S<\theta \ \text{ or } \ r(\hat x_1,s_0)>\lambda,\\
\hat x_1 & \text{otherwise},
\end{cases}
\label{eq:policy}
\end{equation}
\eqentry{Payload routing rule}
and the student finishes with $4-k/2$ steps from this payload.
\begin{whereterms}[2.6cm]
$\pi_{k,\lambda}(x)$ & payload handed to the student for input $x$ \\
$\hat S$, $\theta$ & prompt SNR estimate and fixed gate \\
$r(\hat x_1,s_0)$ & invariant residual (\Cref{eq:residual}) \\
$\lambda$ & residual threshold; $\hat\lambda_k$ after calibration \\
$k$ & teacher path, one of $\mathcal K$
\end{whereterms}
\end{definition}

That is, the router uses the corrected payload when the prompt is clearly noisy or when the teacher's values have moved too far from the draft's, and trusts the teacher otherwise. Algorithm~\ref{alg:handover} states the router as it runs during generation.

\begin{algorithm}[!htbp]
\caption{Risk-calibrated payload router during generation}
\label{alg:handover}
\small
\begin{algorithmic}[1]
\Require prompt mel and mask $M$, text, duration, seed, path $k$, calibrated $\hat\lambda_k$ (or \textsc{none}), gate $\theta$
\If{$\hat\lambda_k=\textsc{none}$} \State \Return output of Algorithm~\ref{alg:baseline} \Comment{fallback} \EndIf
\State $\hat S \gets$ \Cref{eq:snr} on the prompt waveform \Comment{no pass}
\State draw $x_{t_0}\sim\mathcal{N}(0,I)$; run $k$ guided Euler steps on the 8-interval grid \Comment{$2k$ passes}
\State $\hat x_1 \gets x_{t_k} + (1-t_k)\,v(x_{t_k},t_k)$ \Comment{2 passes}
\State $s_0 \gets \varepsilon_0 + g(\varepsilon_0,0)$ \Comment{1 pass}
\State $Q \gets$ \Cref{eq:quantile}; \quad $r \gets$ \Cref{eq:residual} \Comment{no pass}
\If{$\hat S<\theta$ \textbf{or} $r>\hat\lambda_k$} \State $x \gets Q$ \Else \State $x \gets \hat x_1$ \EndIf
\For{$\tau \in \{\tau_{k/2},\dots,\tau_3\}$} \Comment{student, $j=k/2$, $4-k/2$ steps}
  \State $\phi \gets (1-\tau)\,\varepsilon + \tau\,x$; \quad $x \gets \phi + (1-\tau)\,g(\phi,\tau)$ \Comment{1 pass each}
\EndFor
\State remove prompt frames; decode $x$ with the vocoder
\end{algorithmic}
\end{algorithm}

The router has four properties that follow from its definition. On a given path, it always costs the same. Raising the threshold can only send more inputs to the teacher payload. If calibration fails, the released system is used and nothing gets worse. No rule choosing between the two payloads can do better than choosing with hindsight.

\begin{proposition}[Fixed cost and no added network]\label{prop:zero}
For every input and every threshold, the router with path $k$ costs $C(k)$ forward passes, and the SNR estimate, the projection, the residual and the routing decision run no network.
\end{proposition}
\begin{proof}
The teacher path, the draft and the student finish do not depend on the routing decision, so Proposition~\ref{prop:cost} gives $C(k)$ for both payloads. \Cref{eq:snr,eq:quantile,eq:residual,eq:policy} use only frame energies, sorting and norms. Sorting $B$ bands of $n$ frames takes on the order of $Bn\log n$ operations, which is small next to one forward pass of the teacher or student network over the same frames.
\end{proof}

\begin{proposition}[Monotone routing]\label{prop:mono}
For a fixed path $k$, if $\lambda\le\lambda'$, every input that receives the teacher payload under $\pi_{k,\lambda}$ also receives it under $\pi_{k,\lambda'}$; hence $\rho(k,\lambda')\le\rho(k,\lambda)$.
\end{proposition}
\begin{proof}
Under $\pi_{k,\lambda}$ an input receives $\hat x_1$ exactly when $\hat S\ge\theta$ and $r\le\lambda$. Neither $\hat S$ nor $r$ depends on $\lambda$, so $r\le\lambda\le\lambda'$ gives the same payload under $\pi_{k,\lambda'}$.
\end{proof}

\begin{proposition}[Safe fallback]\label{prop:fallback}
If the calibration accepts no threshold for a path and the released system is used instead, the system in use has risk zero.
\end{proposition}
\begin{proof}
If the deployed output is $y_{\text{ref}}(x)$, then $e=F(x,y_{\text{ref}})(1-F(x,y_{\text{ref}}))=0$ for every input.
\end{proof}

\begin{proposition}[Hindsight bound]\label{prop:oracle}
On a given path $k$, let the hindsight router choose, for each input, the teacher payload unless the teacher payload's output fails while the projection's output passes. No rule that chooses between the two payloads of that path has a lower risk.
\end{proposition}
\begin{proof}
For each input the excess-failure indicator of a two-payload rule is the indicator of the chosen payload's output, multiplied by the same factor $1-F(x,y_{\text{ref}})$. The hindsight router picks a failing payload only when both fail, so its indicator is the smallest possible for that input; averaging over inputs preserves the inequality.
\end{proof}
The hindsight router needs to know the outcome of both payloads in advance, so it cannot be used in practice; it will be reported, for each path, as the lower limit of risk that any two-payload rule could reach, next to the risks of always and never projecting.

% =====================================================================
\subsection{Calibration with a Reliability Guarantee}
\label{sec:obj3}

\subsubsection{Calibration Procedure}
\label{sec:calibration}

This section covers Steps~8 to~10 of Table~\ref{tab:steps}. It will set the threshold of every candidate path so that the limit on risk of \Cref{eq:objective} holds, with a stated level of confidence, on new inputs. Calibration uses the Learn-Then-Test procedure \parencite{angelopoulos2025ltt}, which treats each candidate setting as a claim to be tested on data and accepts only settings that the data show to be safe. Before calibration, the failure rule (Definition~\ref{def:failure}, including the minimum levels from the development speakers), the tolerance $\alpha$, the confidence parameter $\delta$, the candidate paths $\mathcal K=\{2,4,6\}$ of \Cref{eq:paths} and the list of candidate thresholds
$\Lambda=\{\lambda_1<\dots<\lambda_G\}$, evenly spaced from zero to one, will be fixed. For every path, the list is ordered from the threshold that sends the most inputs to the projection to the one that sends the fewest, so that the thresholds tested first are the safest; at $\lambda_1=0$ every input with a nonzero residual is projected. The residual is a relative distance, so a value above one would mean that the teacher's estimate differs from its projection by more than its own size; the list therefore covers the whole useful range in small, even steps, and the data rather than the list will limit the chosen threshold. Because the thresholds of a path are tested in a fixed order, adding more candidate thresholds does not weaken the guarantee.

The outcomes of both payloads of every path and of the released system will be recorded for every calibration input with Algorithm~\ref{alg:grid}, each speaker contributing the input of its assigned prompt condition (Section~\ref{sec:data}). The router's output for any path and threshold follows from the recorded $\hat S$ and $r^{(k)}$, so the excess failures of every candidate can be counted without generating speech again. For path $k$ and threshold $\lambda_g$, the observed risk on the $n$ calibration inputs is the share of excess failures:
\begin{equation}
\hat R_{k,g} = \frac1n\sum_{l=1}^n e(x_l,k,\lambda_g).
\label{eq:emprisk}
\end{equation}
\eqentry{Empirical risk on the calibration set}
\begin{whereterms}[2.6cm]
$\hat R_{k,g}$ & observed excess-failure rate of path $k$ with threshold $\lambda_g$ \\
$n$ & number of calibration inputs (one per speaker) \\
$x_l$ & $l$-th calibration input \\
$e(\cdot,\cdot,\cdot)$ & excess-failure indicator (\Cref{eq:excess})
\end{whereterms}

For each candidate, the test starts from the cautious assumption that it is unsafe, $H_{k,g}: R(k,\lambda_g)>\alpha$, and accepts the candidate only if the data give strong evidence against that assumption. The evidence is measured by a p-value from the binomial distribution, which describes how many failures to expect among $n$ independent inputs:
\begin{equation}
p_{k,g} = \P\big(\mathrm{Bin}(n,\alpha)\le n\hat R_{k,g}\big),
\label{eq:pvalue}
\end{equation}
\eqentry{Binomial p-value}
\begin{whereterms}[2.6cm]
$p_{k,g}$ & p-value of path $k$ with threshold $\lambda_g$ \\
$\mathrm{Bin}(n,\alpha)$ & number of excess failures among $n$ independent inputs if the true risk were exactly $\alpha$ \\
$n\hat R_{k,g}$ & observed number of excess failures of that candidate
\end{whereterms}
that is, the probability of seeing so few excess failures if the true risk were exactly $\alpha$. A small p-value means that so few failures would be unlikely for an unsafe candidate. Because three paths are calibrated and later compared, the allowance $\delta$ is divided equally among them: on each path, a threshold is accepted as safe when $p_{k,g}\le\delta/\lvert\mathcal K\rvert=\delta/3$. This keeps the guarantee valid for all three paths at once, so it still holds for whichever path the evaluation later finds best. The thresholds of each path are tested in the order of the list, and the test of that path stops at the first threshold that is not accepted, returning $\hat\lambda_k$, the last accepted threshold, or nothing if even the first is not accepted; in that case the path falls back to the released system (Proposition~\ref{prop:fallback}). Algorithm~\ref{alg:ltt} states the procedure.

\begin{algorithm}[!htbp]
\caption{Learn-Then-Test calibration of the residual threshold of each path}
\label{alg:ltt}
\small
\begin{algorithmic}[1]
\Require records of Algorithm~\ref{alg:grid} for $n$ calibration inputs; paths $\mathcal K=\{2,4,6\}$; grid $\lambda_1<\dots<\lambda_G$; $\alpha$, $\delta$
\For{$k \in \mathcal K$} \Comment{each path tested separately at level $\delta/3$}
  \State $\hat\lambda_k \gets \textsc{none}$
  \For{$g = 1,\dots,G$}
    \State $\hat R_{k,g} \gets \frac1n\sum_{l} e(x_l,k,\lambda_g)$ \Comment{from records, no new generation}
    \State $p_{k,g} \gets \P\big(\mathrm{Bin}(n,\alpha)\le n\hat R_{k,g}\big)$
    \If{$p_{k,g}\le\delta/\lvert\mathcal K\rvert$}
      \State $\hat\lambda_k\gets\lambda_g$
    \Else
      \State \textbf{break}
    \EndIf
  \EndFor
\EndFor
\State \Return $\hat\lambda_k$ for every $k\in\mathcal K$ \Comment{\textsc{none}: path uses the released system}
\end{algorithmic}
\end{algorithm}

\subsubsection{Guarantee}

The procedure can be trusted for three reasons. Each p-value is small only rarely when its candidate is in fact unsafe. Within a path, testing in a fixed order means that the test can go wrong only by accepting the first unsafe threshold in the list, so that path's whole share $\delta/3$ protects that one test instead of being divided among all thresholds. Across paths, the three shares add up to $\delta$, so the chance that any path accepts an unsafe threshold is at most $\delta$. Learn-Then-Test needs only a pass-or-fail outcome per input, calibration inputs that are independent and come from the same population as future inputs, and candidates fixed in advance \parencite{angelopoulos2025ltt}; the excess-failure indicator of \Cref{eq:excess} is such an outcome, so the procedure applies to text-to-speech without change.

\begin{proposition}[Risk control]\label{prop:ltt}
If the calibration inputs and a new input are independent draws from $\mathcal{D}$, then
\begin{equation}
\P\big(R(k,\hat\lambda_k)\le\alpha\ \text{for every}\ k\in\mathcal K\big)\ \ge\ 1-\delta,
\label{eq:guarantee}
\end{equation}
\eqentry{Reliability guarantee}
where the risk of the fallback is zero. Put simply, with confidence at least $1-\delta$, every calibrated path adds failures on at most a share $\alpha$ of new inputs; by Proposition~\ref{prop:rate}, the router then fails on at most a share $\alpha$ more of the new inputs than the released system.
\begin{whereterms}[2.6cm]
$R(k,\hat\lambda_k)$ & true excess-failure rate of the router on path $k$ on new inputs \\
$\hat\lambda_k$ & threshold of path $k$ returned by Algorithm~\ref{alg:ltt} \\
$1-\delta$ & confidence; the probability refers to the random choice of calibration speakers
\end{whereterms}
\end{proposition}
\begin{proof}
\emph{Valid p-values.} If $H_{k,g}$ is true, $n\hat R_{k,g}\sim\mathrm{Bin}(n,R(k,\lambda_g))$ with $R(k,\lambda_g)>\alpha$. A binomial count with a larger rate tends to be larger, so a small count is less likely than under $\mathrm{Bin}(n,\alpha)$; hence $\P(p_{k,g}\le\delta/3)\le\delta/3$.
\emph{One path.} Let $A_k$ be the event that path $k$ returns an unsafe threshold, $R(k,\hat\lambda_k)>\alpha$. Then a true hypothesis $H_{k,\hat g}$ was rejected. The test of path $k$ rejects every hypothesis before $\hat g$, so the first true hypothesis of that path in the fixed order, $g^\ast_k$, which depends only on the unknown risks and not on the data, was rejected. Hence $A_k$ implies $p_{k,g^\ast_k}\le\delta/3$, and $\P(A_k)\le\delta/3$. This is the fixed-sequence argument of \textcite{bauer1991} as used by \textcite{angelopoulos2025ltt}.
\emph{All paths.} The guarantee fails only if some $A_k$ happens. The chance that at least one of several events happens is at most the sum of their chances, so the probability of failure is at most $\sum_{k\in\mathcal K}\P(A_k)\le 3\cdot\delta/3=\delta$. Returning the fallback never errs (Proposition~\ref{prop:fallback}).
\end{proof}

The proof does not need the risk to change in any particular way with $\lambda$ or $k$. The guarantee is an average over inputs and over the mix of prompt conditions, not a promise about a single sentence or a single condition, and it does not hold if new inputs differ systematically from the calibration inputs.

\subsubsection{Tolerance, Confidence, and Calibration Size}
\label{sec:calsize}

The number of calibration speakers determines how small a risk can be shown at all. With no excess failures among $n$ calibration inputs, the p-value of \Cref{eq:pvalue} is $(1-\alpha)^n$ \parencite{clopper1934}, so a threshold can be accepted on a path only if
\begin{equation}
(1-\alpha)^n\le\frac{\delta}{\lvert\mathcal K\rvert}, \qquad\text{that is,}\qquad n \ge \frac{\ln(\delta/\lvert\mathcal K\rvert)}{\ln(1-\alpha)}.
\label{eq:zerofail}
\end{equation}
\eqentry{Minimum calibration size with zero failures}
\begin{whereterms}[2.6cm]
$n$ & number of independent calibration inputs (one per speaker) \\
$\alpha$ & tolerated risk \\
$\delta$ & allowed probability that calibration accepts an unsafe threshold on some path \\
$\lvert\mathcal K\rvert$ & number of candidate paths, 3 \\
$\ln$ & natural logarithm
\end{whereterms}
\Cref{eq:zerofail} links the three quantities that the user of the method chooses: a smaller tolerance or a higher confidence needs more calibration speakers, and testing more candidate paths needs more as well, because each path receives a smaller share of $\delta$. The tolerance and the confidence are design choices; Learn-Then-Test leaves both to the user \parencite{angelopoulos2025ltt}. The values of $\alpha$, $\delta$ and the other settings of the study will be fixed in the study protocol before calibration and reported with the results; Appendix~\ref{app:calc-size} shows the calculation for one choice of values. Because the choice of $\alpha$ is a judgment, the calibration will also be repeated at a larger tolerance as a sensitivity analysis.

\subsubsection{Verifying the Guarantee and Reporting the Thresholds}
\label{sec:verify}

This check begins the evaluation (Step~11). The guarantee is a statement about repeated calibration, so it will be checked in two ways. The single check on new speakers will apply each path with its threshold $\hat\lambda_k$ to the test split, each speaker with its assigned condition, and report the observed excess-failure rate with a Clopper--Pearson interval \parencite{clopper1934}; it will also report the failure rates of the router and of the released system on the same inputs, so that \Cref{eq:rate} can be seen directly. A rate well above $\alpha$ would indicate that an assumption does not hold, most likely that the inputs are independent or that calibration and test inputs come from the same population. The repeated check will combine the calibration and test speakers, divide them at random into two halves $R$ times, run Algorithm~\ref{alg:ltt} on one half and measure the risk of every path on the other each time. The guarantee predicts that the risks of all paths stay within $\alpha$ in at least a share $1-\delta$ of these divisions. Because all outputs will already be recorded, neither check requires generating speech again.

The calibrated threshold $\hat\lambda_k$ of every path will be reported so that other researchers can reuse or reproduce it, together with the conditions under which it holds: $\alpha$, $\delta$, $n$, the candidate paths and their testing order, the grid, the student's starting point, the SNR gate, the constants and minimum levels of the failure rule, the trained networks, the evaluation models, the dataset and the mixture of prompt conditions, and the hardware. Under these conditions, the thresholds can be used directly with DMOSpeech~2. For another model, dataset, risk limit or failure rule they do not transfer, but the procedure does: Algorithm~\ref{alg:grid} is run on enough independent calibration inputs to satisfy \Cref{eq:zerofail}, and Algorithm~\ref{alg:ltt} returns new thresholds with the same kind of guarantee. Calibrated thresholds of this kind are established in other settings: CALM certified when a language model may stop computing early \parencite{schuster2022calm}, thought calibration used Learn-Then-Test with fixed-sequence testing to certify when a model may stop reasoning \parencite{wu2025thought}, and FlowGuard calibrated a threshold on a flow-matching sampler \parencite{li2025flowguard}. What these studies did not provide is such a threshold for the teacher-to-student handover in text-to-speech, with a bound stated against the released system's own failure rate.

% =====================================================================
\subsection{Evaluation and Robustness}
\label{sec:obj4}
\label{sec:analysis}

\subsubsection{Compared Systems}
\label{sec:compared}

The evaluation covers Steps~12 to~14 of Table~\ref{tab:steps}. All systems in Table~\ref{tab:systems} will be run on the test split under all eight prompt conditions and on LibriSpeech-PC with clean prompts. The established systems show what the released models achieve: the released DMOSpeech~2 system is the baseline that the router must not fall below, and the teacher alone and the student alone mark the expensive and the cheap end of the range. The proposed system is the router on each candidate path with its calibrated threshold. The three paths give a curve of cost against quality, variety and excess failures; the least costly path whose excess-failure rate stays within the limit and whose quality is not worse than the released system's will be identified in the conclusions. The ablations, versions of the router with one part removed or fixed, will be run on every path and show what each part contributes: never and always projecting show what each payload does without a decision; the SNR-only router shows what the residual adds beyond the prompt check; and the student restarted from its own draft shows what the teacher adds at all. The hindsight router (Proposition~\ref{prop:oracle}) will be reported as a limit, not as a competitor. The results will be reported in Chapter~4 as a comparison matrix against the baseline, which reports every system of Table~\ref{tab:systems} on the measures of Table~\ref{tab:matrix-cols}; the matrix will be given overall, for every prompt condition, and for LibriSpeech-PC.

\begin{table}[!htbp]
\centering\small
\caption{Compared systems}
\label{tab:systems}
\label{tab:methods}
\begin{tabularx}{\textwidth}{@{}>{\raggedright\arraybackslash}p{4.6cm} >{\raggedright\arraybackslash}X >{\centering\arraybackslash}p{1.4cm}@{}}
\toprule
\textbf{System} & \textbf{What it tests} & \textbf{Passes} \\
\midrule
\multicolumn{3}{@{}l}{\textit{Established systems}} \\
DMOSpeech~2, released system ($s=0.25$) & Baseline and reference of the risk & 21 \\
F5-TTS teacher alone & Quality of the teacher without distillation & 64 \\
DMOSpeech~2 student alone & Cheapest released system & 4 \\
\midrule
\multicolumn{3}{@{}l}{\textit{Proposed}} \\
Payload router \texttt{SNR\_R}, each path $k$ with $\hat\lambda_k$ & The method; cost against quality and risk across paths & $C(k)$ \\
\midrule
\multicolumn{3}{@{}l}{\textit{Ablations (each path $k$)}} \\
Always project ($Q$ for every input) & The projection without a decision & $C(k)$ \\
Never project ($\hat x_1$ for every input) & The teacher payload without a decision; the effect of the shorter teacher path alone & $C(k)-1$ \\
SNR-only router ($Q$ if $\hat S<\theta$) & Value of the residual beyond the prompt check & $C(k)-1$ or $C(k)$ \\
Student from its draft (\texttt{STUR2}) & Value of the teacher at all & 3 \\
\midrule
Hindsight router & Lowest risk of any two-payload rule (descriptive) & --- \\
\bottomrule
\end{tabularx}
\end{table}




\subsubsection{Measures of Comparison}
The saving of a system is stated relative to the released system:
\begin{equation}
\Delta C = 1-\frac{\bar C_{\text{sys}}}{C_{\text{ref}}}.
\label{eq:saving}
\end{equation}
\eqentry{Relative saving in forward passes}
\begin{whereterms}[2.6cm]
$\Delta C$ & share of the reference's forward passes saved \\
$\bar C_{\text{sys}}$ & mean forward passes per utterance of the system \\
$C_{\text{ref}}$ & forward passes of the released system, 21
\end{whereterms}
The same ratio will be reported for measured generation time on the Tesla T4. For the router, the routing rate $\rho(k,\hat\lambda_k)$ of each path will be reported overall and per condition, together with the share of clean prompts that receive the projection, since projecting a clean prompt gives up teacher detail for no gain in safety.

\subsubsection{Robustness}
\label{sec:robust}

Three robustness questions follow from the assumptions of Proposition~\ref{prop:ltt}. The first concerns prompt conditions: the excess-failure rate and the quality measures will be reported for each of the eight conditions, and the thresholds will also be calibrated on clean prompts only and tested on each damaged condition, to see how the guarantee and the residual behave when the prompts met in use differ from those used in calibration. The second concerns the number of calibration speakers: the calibration will be repeated on random subsets of calibration speakers of several sizes, $R$ times each, and the study will report how often each path's threshold is accepted, the resulting routing rate and the observed risk, which shows how much data a user of the method needs. The third concerns the tolerance: the calibration will be repeated at a larger tolerance, as stated in Section~\ref{sec:calsize}. Transfer to a new source will be checked on LibriSpeech-PC with the thresholds calibrated on Seed-TTS.

\subsubsection{Statistical Analysis}
\label{sec:stats}

Comparisons will be paired, since every system is applied to the same inputs with the same random starting values; each speaker's result under one system is compared with the same speaker's result under another. Differences in generation time, WER, SIM, UTMOS and $\mathrm{CV}_{f_0}$ between the router and each other system will be tested with the Wilcoxon signed-rank test \parencite{wilcoxon1945}, which compares paired values without assuming that the differences follow a normal distribution. Because the router is compared with several systems, the Holm correction \parencite{holm1979} will adjust the results so that the chance of at least one false finding stays at the chosen level. The teacher passes of the router are justified only by the variety they add, so one comparison is fixed in advance as the main test of the teacher's value: the router's $\mathrm{CV}_{f_0}$, raw and aligned, against that of the student restarted from its own draft, with a one-sided Wilcoxon signed-rank test on speaker-level values. Because pitch-based variety captures only part of what listeners hear \parencite{yang2025prosody}, a small listening check by the research team on a fixed subset of these outputs will be reported as supporting evidence only. Whether two systems fail on different inputs will be compared with the exact McNemar test \parencite{mcnemar1947}, which looks only at the inputs on which one system fails and the other does not. Average differences will be reported with bootstrap intervals, obtained by resampling whole speakers $B$ times \parencite{efron1979bootstrap}; resampling speakers rather than single sentences respects the dependence between sentences of the same speaker \parencite{liu2020blockwise}.

A claim that the router's quality is not worse than the released system's will require a test of equivalence, the two one-sided tests procedure \parencite{schuirmann1987}, which shows that the difference lies within a stated margin; failing to find a difference is not the same as showing that there is none, so it will not be taken as evidence of equal quality. The margins $\Delta_{\text{WER}}$, $\Delta_{\text{SIM}}$ and $\Delta_{\text{UTMOS}}$, the largest losses against the released system's mean that still count as ``not worse'', will be fixed before testing. Results will be reported per condition as well as overall, so that a gain on damaged prompts cannot hide a loss on clean ones.

Every test will use a significance level $a$ after the Holm correction. This level is a different quantity from the tolerance $\alpha$ of the calibration, even if the two values are equal: the significance level limits the chance of declaring a difference between two systems when there is none, while $\alpha$ limits how often the router may fail where the released system does not.

% =====================================================================
\subsection{Ethical Considerations}
\label{sec:ethics}

Voice cloning can be used to imitate real people without their consent. The study will therefore use only public research datasets under their licenses \parencite{ardila2020commonvoice,panayotov2015librispeech}, will not release cloned voices beyond what those licenses permit, and will report aggregate results rather than distribute synthesized speech of identifiable speakers. It will not involve end users and will make no claims about effects on autistic users or other vulnerable groups; any small listening check will be limited to the research team.
